\section{Where a deleted matching sends the zeros}\label{sec:zeros}
Let $G_{d,M}$ be $K_{d+1}$ with a matching of $M$ edges deleted, with
$0\le M\le\lfloor(d+1)/2\rfloor$ and $d\ge2$. Write $E_{d,M}$ for
its Ehrhart polynomial, $H_d=E_{d,0}$, and $C_d$ for the Ehrhart polynomial
of the $d$-dimensional cross-polytope. An omitted edge is a part of size
two. In \eqref{eq:multipartite-gamma} such a part changes only
$\gamma_1$, by $-2$. Since $4t=(1+t)^2-(1-t)^2$, comparison of the
Ehrhart series gives the exact identity
\begin{equation}\label{eq:matching}
 E_{d,M}(s)=H_d(s)-\tfrac M2\big(C_d(s)-C_{d-2}(s)\big).
\end{equation}
Here and below $s$ is the complex dilation variable. The palindromic
numerators give $E_{d,M}(-1-s)=(-1)^d E_{d,M}(s)$, so conjugation
and reflection across $\Re s=-1/2$ both preserve the zero multiset.

\begin{theorem}[Bounded-height zero counts]\label{thm:strings}
Fix $0<t_1<t_2<\infty$, write $I=[t_1,t_2]$, $L=\log d$, and let
$d\to\infty$.
\begin{enumerate}
\item If $M=o(\sqrt d)$, then in every fixed rectangle
$|\Re s+1/2|\le R$, $\Im s\in I$, $R>0$, all the zeros are simple
and on $\Re s=-1/2$ for sufficiently large $d$. Their number is
$2|I|L/\pi+O(1)$, and consecutive spacings are
$(\pi/2+o(1))/L$.
\item Suppose $M/d^\alpha\to\lambda>0$, $1/2<\alpha\le1$,
subject to the matching-size bound. In any fixed rectangle as above
containing the following three real parts in its interior, the zeros with
height in $I$ are simple for sufficiently large $d$. Their locations
and counts are
\[
\begin{array}{c|c|c}
 \text{real part} & \text{number} & \text{spacing in height}\\ \hline
 -\alpha+O(1/L) & |I|L/(2\pi)+O(1)&(2\pi+o(1))/L\\
 -1/2\ \text{exactly} & |I|L/\pi+O(1)&(\pi+o(1))/L\\
 \alpha-1+O(1/L) & |I|L/(2\pi)+O(1)&(2\pi+o(1))/L .
\end{array}
\]
Spacing is measured within a single row. For a root
$s=\sigma+it$ in the right-hand row one has the sharper formula
\begin{equation}\label{eq:shift}
 L(\sigma-\alpha+1)
 =\log\left|\lambda\,2^{\alpha-1+it}(\alpha-1+it)
               \Gamma(\alpha+it)\right|+o(1).
\end{equation}
\end{enumerate}
All errors in the location and spacing statements are uniform on $I$.
\end{theorem}

For example, deleting $\lfloor d^{3/4}\rfloor$ pairwise disjoint edges
splits the bounded-height zeros between the lines with real parts
$-3/4,-1/2,-1/4$. A largest matching gives limiting real parts $-1,-1/2,0$.
In all of these graphs the fraction of edges deleted tends to zero.
The proportions in Theorem~\ref{thm:strings} concern bounded heights,
where only order $\log d$ zeros in a bounded rectangle are counted.

\subsection{Two dimension generating functions}
Define $s^*=-1-s$, $\epsilon=(-1)^d$, and
\[
 A(s)=\Gamma(1+s)^{-2},\qquad B(s)=2^s/\Gamma(s).
\]
\begin{lemma}\label{lem:expansion}
Uniformly on compact sets of $s=\sigma+it$ with $t\ne0$, one has
\begin{align}
 H_d(s)&=A(s)d^{2s}+\epsilon A(s^*)d^{2s^*}
       +O\big(d^{2\sigma-1}+d^{-3-2\sigma}\big),\label{eq:Hexp}\\
 \tfrac12(C_d(s)-C_{d-2}(s))
   &=B(s)d^{s-1}+\epsilon B(s^*)d^{s^*-1}
       +O\big(d^{\sigma-2}+d^{-3-\sigma}\big).\label{eq:Cexp}
\end{align}
The same estimates hold on compact sets avoiding the resonances
$1+2s\in\mathbb Z$.
\end{lemma}
\begin{proof}
Use the dimension variable $z$. The cross-polytope counts lattice points
by their support, so $C_d(s)=\sum_j 2^j\binom dj\binom sj$. Summing
yields
\[
 \sum_{d\ge0} C_d(s)z^d=(1+z)^s(1-z)^{-1-s}.
\]
For the complete graph, Theorem~\ref{thm:gamma} gives
$h^*(x)=\sum_{j=0}^d\binom dj^2 x^j$. Indeed the coefficient of $x^j$
can be counted by two $j$-subsets according to their symmetric
difference. Its dimension generating function is
\begin{equation}\label{eq:dimension}
 \sum_{d\ge0}H_d(s)z^d=(1-z)^{-1-2s}
       {}_2F_1(-s,-s;1;z^2).
\end{equation}
For details, set $H_0=1$. The elementary identity
\[
 \sum_{d\ge0} u^d\sum_j\binom dj^2 x^j
 =\big(1-2(1+x)u+(1-x)^2u^2\big)^{-1/2}
\]
follows by taking the constant term in $y$ of
$[1-u(1+y)(1+x/y)]^{-1}$. The bivariate generating function
$\sum_{n,d\ge0}H_d(n)x^nz^d$ is therefore
$[(1-z)^2-2(1+z^2)x+(1+z)^2x^2]^{-1/2}$.
Applying the same elementary identity with $u=x/(1-z)^2$ and
argument $z^2$ gives \eqref{eq:dimension} at nonnegative integers $s=n$.
For fixed $d$, its $z^d$ coefficient is a polynomial in $s$ of degree
at most $d$, proving the identity for complex $s$.

At $z=1$ the connection formula for the hypergeometric function at argument one
\cite[Eq.~15.8.4]{dlmf} writes the right side of
\eqref{eq:dimension} as a function analytic near $1$ plus
\[
 (1-z)^{-1-2s}\frac{\Gamma(1+2s)}{\Gamma(1+s)^2}
 (1+O(1-z)).
\]
At $z=-1$ the corresponding singular part is
$(1+z)^{1+2s}\Gamma(-1-2s)/\Gamma(-s)^2$ times
$1+O(1+z)$. The cancellation of analytic factors here uses
$(1-z^2)^{1+2s}(1-z)^{-1-2s}=(1+z)^{1+2s}$.
For the second line of the lemma use instead
$\tfrac12(1+z)^{s+1}(1-z)^{-s}$, obtained by multiplying the
cross-polytope series by $(1-z^2)/2$.

These functions continue analytically to a disk of radius $1+\eta$
cut outward at $1$ and $-1$. In the hypergeometric case take the
preimage of its standard slit under $z\mapsto z^2$. Transfer at the two
singularities now gives the stated terms, using the binomial transfer
coefficient $d^{w-1}/\Gamma(w)$ of $(1-z)^{-w}$.
Specifically the first omitted singular power at each point supplies
one additional factor $d^{-1}$. The analytic local terms can be
subtracted to arbitrarily high order by a polynomial matching their
jets at both points. The indented contours can be chosen the same for
parameters in the given compact set, and the convergent local
connection series have uniformly bounded remainders there. This gives
the uniform errors by multiple-singularity transfer
\cite[Ch.~VI]{fs}.
\end{proof}

\subsection{Counting and locating the zeros}
Put $w=\exp((s+1/2)L)$ and $\mu=M/\sqrt d$. Combining the two expansions
with \eqref{eq:matching} gives
\begin{align}
 dE_{d,M}(s)={}&A(s)w^2+\epsilon A(s^*)w^{-2}
        -\mu\{B(s)w+\epsilon B(s^*)w^{-1}\}+\mathcal E,\label{eq:four}\\
 |\mathcal E|\ll{}&d^{-1}
 \big(|w|^2+|w|^{-2}+\mu(|w|+|w|^{-1})\big).
 \label{eq:error}
\end{align}
On a compact rectangle of positive heights all four coefficients
are bounded and bounded away from zero.

\begin{proof}[Proof of Theorem~\ref{thm:strings}]
If $\mu\to0$, all zeros satisfy $|\log|w||\le C$ for a fixed
$C$. Outside this band one of $w^2,w^{-2}$ dominates the other three
terms and the error. Inside, divide \eqref{eq:four} by $A(s)w^{-2}$.
The resulting equation for a zero is
\[
 w^4=-\epsilon A(s^*)/A(s)+o(1).
\]
Here and in what follows the errors hold for holomorphic functions
uniformly on bands of width $C/L$, for arbitrarily large fixed $C$.

If $\mu\to\infty$ in part (2), set $x=|w|$. In order of increasing
$x$, the controlling magnitudes are
$x^{-2},\mu x^{-1},\mu x,x^2$.
Their consecutive crossings are $x=\mu^{-1},1,\mu$.
Outside a fixed multiplicative neighborhood of each crossing, a
single term dominates. Choose the neighborhood factor larger than
ten times the ratio of the maximum to the minimum coefficient
modulus on a slightly enlarged rectangle. Equation~\eqref{eq:error}
then excludes zeros there. Near $x=1$, division by the central terms
gives $w^2=-\epsilon B(s^*)/B(s)+o(1)$. Near $x=\mu$, division
by $A(s)\mu w$ gives
\[
  e^{(s-\alpha+1)L}
    = \frac{M}{d^\alpha}\,\frac{B(s)}{A(s)}+o(1).
\]
Reflection gives the left band. Taking absolute values on the right
gives \eqref{eq:shift} by the gamma recurrence.

For clarity, each reduced equation has the form
$e^{k(s-a)L}=g_d(s)+o(1)$ on $|\Re s-a|\le C/L$.
Here $k=4,2,1$ in the three cases (subcritical, central supercritical,
right supercritical). At each fixed parity $g_d$ is a nonvanishing
analytic function times a positive scalar tending to a positive limit.
It has an analytic logarithm $\ell_d$ with bounded value and first
derivative on the enlarged rectangle. For each integer $j$ whose height
$y_j=2\pi j/(kL)$ is in a slightly larger interval, the leading equation
has a unique nearby root
\[
 s_j=a+iy_j+\frac{\ell_d(a+iy_j)}{kL}+O(L^{-2}).
\]
Indeed $s\mapsto a+iy_j+\ell_d(s)/(kL)$ is a contraction on a disk
of radius $C'/L$ about $a+iy_j$. On a circle about $s_j$ of radius
$\eta_d/L$, the leading difference has modulus $\gg\eta_d$ for
$\eta_d\to0$ slowly enough to dominate the uniform error. Rouch\'e's
theorem gives one simple true zero in each disk. Conversely, taking
a logarithm at a true zero yields $k(s-a)L=2\pi i j+\ell_d(s)+o(1)$,
placing it in one of the disks, enlarging $\eta_d$ if necessary.
This proves the counts, spacings, and simplicity, including the
$O(1)$ possible endpoint adjustments. For the central equations
$|g_d(-1/2+it)|=1$. The leading roots there lie on the canonical line
by contraction restricted to that line. Their disks
are invariant under $s\mapsto-1-\bar s$. Uniqueness puts the true
zero exactly on that line.
\end{proof}

\begin{corollary}[The entire limiting zero set]\label{cor:strip}
Over all admissible choices of $M=M(d)$, the set of finite limits of
zeros of $E_{d,M}$ along sequences $d\to\infty$ is exactly
\[
                     \{s\in\mathbb C:-1\le\Re s\le0\}.
\]
\end{corollary}
\begin{proof}
For every real part in $(-1/2,0]$ choose $\alpha=1+\Re s$ and a
permitted positive $\lambda$ in Theorem~\ref{thm:strings}(2).
Its spacing assertion supplies a sequence of zeros converging to
every point of that line of positive height. Conjugation, reflection,
and part (1) supply the remaining nonreal points. A diagonal sequence
with heights tending to zero gives every point of the real interval.

Conversely, on each compact subset of $\Re s>0$ avoiding the real
resonances, \eqref{eq:Hexp}, \eqref{eq:Cexp} and $M=O(d)$ give
$d^{-2s}E_{d,M}(s)\to A(s)$ uniformly over all allowed matchings.
The same convergence holds near each resonance by applying the maximum
modulus principle to the difference on a small circle contained in
that half-plane and avoiding all resonances. The limit $A$ has no zeros
in the half-plane. Reflection excludes limits to the left of $-1$.
\end{proof}

\section{The threshold at the square-root scale}\label{sec:critical}
A finer question asks whether deleting just $c\sqrt d$ edges can
already produce zeros off the canonical line. The threshold has a
closed form if one asks for uniform stability over both dimension
parities. Below, a bounded region can be taken to be any fixed compact
rectangle in the complex plane.

\begin{theorem}[Sharp parity-uniform threshold]\label{thm:critical}
Assume $M/\sqrt d\to c\ge0$ and set $c_*= \sqrt{8/\pi}$.
\begin{enumerate}
\item If $c<c_*$, all zeros in any bounded region are simple and lie on
the canonical line for sufficiently large $d$.
\item If $c>c_*$, all sufficiently large even $d$ have off-line zeros
converging to $-1/2$. More precisely, set $r=c\sqrt{\pi/8}$,
$x_-=(-r-\sqrt{r^2+8})/2$, and
$h=\operatorname{arcosh}(-x_-/2)>0$. There are zeros
\[
 -\tfrac12+\frac{\pm h+i\pi}{\log d}
                  +o(1/\log d)
\]
and their complex conjugates. For odd $d$ the same conclusion holds
when $r>2$, with $h=\operatorname{arcosh}(r/2)$.
\end{enumerate}
\end{theorem}

\begin{proof}
At positive and negative heights, \eqref{eq:four} is uniform. It also
has absolute error $o(1)$ on a small fixed disk about $-1/2$ for bounded
$\mu$. Indeed on a circle of radius $0<\delta<1/4$ about that point
\eqref{eq:error} gives $O(d^{-1+2\delta})$. Both sides of
\eqref{eq:four} are holomorphic in $s$, and the maximum modulus principle
extends the absolute bound to the disk. On a compact set strictly
to the right of the canonical line, $d^{-2s}E_{d,M}(s)\to A(s)$
by the four-term expansion. The same small-circle argument for this
normalized difference extends across any real resonance in that open
half-plane. Reflection gives the corresponding exclusion on the left.
Near the canonical line the coefficients $A(s),A(s^*)$ are separated
from zero and all four coefficients are bounded. Dominance at large or
small $|w|$ gives $|\log|w||=O(1)$ for a zero in a bounded region;
the absolute $o(1)$ error on the exceptional disk suffices as well.

Fix one parity $\epsilon$ and any sequence of real heights $t_d\to t$.
On a subsequence $e^{it_d L}\to\tau$, $|\tau|=1$. At
$s=-1/2+it_d+u/L$, \eqref{eq:four} converges locally uniformly in $u$
to $P(\tau e^u)/(\tau e^u)^2$, where, with $s_0=-1/2+it$,
\begin{equation}\label{eq:quartic}
 P(w)= A(s_0) w^4-c B(s_0)w^3
             -\epsilon c\overline{B(s_0)}w+\epsilon\overline{A(s_0)} .
\end{equation}
On the unit circle the equation $P(w)=0$ is
$\Re(\kappa A(s_0)w^2)=c\Re(\kappa B(s_0)w)$, with $\kappa=1$ for
even $d$ and $\kappa=-i$ for odd $d$.
If $c|B(s_0)/A(s_0)|<1$, the signs alternate at the four successive
extrema of the left side. This gives four distinct unit roots of a
quartic, necessarily simple. The gamma reflection identity gives
\[
 \left|\frac{B(-1/2+it)}{A(-1/2+it)}\right|^2
     =\frac{\pi(1/4+t^2)}{2\cosh(\pi t)}
     \le\frac{\pi}{8},
\]
since $\cosh(\pi t)\ge1+\pi^2t^2/2\ge1+4t^2$.
This proves the strict inequality uniformly for $c<c_*$.

To see that exact finite-$d$ roots are on the line, argue by
contradiction along a sequence of multiple or off-line roots and take
$t_d$ to be their imaginary parts. Their rescaled real offsets $u_d$
are bounded and have a convergent subsequence. A unit-modulus limiting
root forces $u_d\to0$. Simplicity in the limiting quartic implies that
a fixed small disk about $u=0$ contains just one zero counted with
multiplicity for all large $d$, by Rouch\'e's theorem. Reflection
$u\mapsto-\bar u$ preserves this disk and the finite zero set, so that
zero is simple and has real part zero. This contradiction proves (1).

At $t_d=t=0$ one has $\tau=1$, $A=1/\pi$ and
$B=-1/(2\sqrt{2\pi})$. For even $d$ the limiting quartic divided
by $A w^2$ is $x^2+r x-2$, $x=w+w^{-1}$. If $r>1$, its lower
root $x_-$ is less than $-2$, giving the two simple roots
$w=-e^{\pm h}$. For odd $d$ the quartic divided by $A$ factors as
$(w^2-1)(w^2+r w+1)$; this gives two simple negative roots off the
unit circle for $r>2$. Rouch\'e at $u=\pm h+i\pi$ proves (2).
At equality the limiting root is multiple; the strict assertions above
do not decide that case.
\end{proof}

\paragraph{The critical law at other heights.}
Equation~\eqref{eq:quartic} also determines the local transition at each
fixed height. For an explicit test rotate $w$ so that the unit-circle
equation becomes $\cos(2\theta)+a\cos\theta+b\sin\theta=0$.
Explicitly set $D=\kappa B(s_0)e^{-i\arg(\kappa A(s_0))/2}/|A(s_0)|$
and $(a,b)=c(-\Re D,\Im D)$. A double root requires the equation and
its derivative to vanish, giving the boundary
\[
 (a,b)=(2\cos^3\theta-3\cos\theta,\,
              3\sin\theta-2\sin^3\theta).
\]
This curve meets every ray once. Its radius squared is
$1+3\sin^2(2\theta)$ and its polar-angle derivative has numerator
$ab'-ba'=-3\cos^2(2\theta)$. Its four quadrants trace one turn, as
can also be seen from the signs of $a$ and $b$. The number of unit
roots is constant on either side of the curve. The origin has four,
and a point sufficiently far along the real axis has two, by the
quadratic in $\cos\theta$. Inside the boundary all four roots are
simple on the unit circle; outside, two remain there and the other
two form a reciprocal-conjugate pair. This description concerns
the quartic at regular parameter values. Multiple roots on the
boundary require higher-order information.
